\documentclass[12pt]{article}
\usepackage{latexblog}

% ============================================================
% Blog Metadata
% ============================================================
\blogtitle{A brief introduction to MySQL binary log}
\blogdate{2020-06-15}
\blogtags{MySQL, 闲话编程, 存储}
\bloglang{en}
\blogsummary{}

\begin{document}
\makeblogtitle

Accroding to the mysql manual, the binary log( also known as binlog) is a very powerful feature that enables the recording of ``events'' that describe database changes. Those changes could be table modification or data change, and may also contain some statements which may potentially made changes such as \texttt{DELETE} which no matched items.

This kind of feature enables mysql in replication from master to slave servers by sending events contained in binlogs, and also data recovery. Except for that, it could also be used in CDC(change-data-capture) since it's event based, an example usage could be found in the cdc-component in \href{https://eventuate.io/}{Eventuate™}.

\section{binary log configuration}\label{binary-log-configuration}

\subsection{enable mysql binlog}\label{enable-mysql-binlog}

By default the binlog feature is not enabled by mysql, we can enable this by adding the following settings to \texttt{my.cnf} file:

\begin{Shaded}
\begin{Highlighting}[]
\KeywordTok{[mysqld]}
\DataTypeTok{log{-}bin}\OtherTok{=}\StringTok{mysql{-}bin}
\DataTypeTok{server\_id}\OtherTok{=}\DecValTok{1}
\end{Highlighting}
\end{Shaded}

Since MySQL5.7, the \texttt{server\_id} must be specified if binlog is enabled, and this value should be unique among the cluster. Then we can find binlog like the following:

\begin{verbatim}
mysql> show binary logs;
+------------------+-----------+
| Log_name         | File_size |
+------------------+-----------+
| mysql-bin.000001 |       154 |
+------------------+-----------+
1 row in set (0.00 sec)
\end{verbatim}

\subsection{binlog format}\label{binlog-format}

There are 3 kinds of binlog format supported by mysql:

\begin{itemize}
\tightlist
\item
  STATEMENT: statement-based logging
\item
  ROW: row-based logging
\item
  MIXED: in mixed mode
\end{itemize}

By default, the statement-based binlog format is used, and that could be changed by settings:

\begin{Shaded}
\begin{Highlighting}[]
\KeywordTok{[mysqld]}
\DataTypeTok{log{-}bin}\OtherTok{=}\StringTok{mysql{-}bin}
\DataTypeTok{server\_id}\OtherTok{=}\DecValTok{1}
\DataTypeTok{binlog\_format}\OtherTok{=}\StringTok{"ROW"}
\end{Highlighting}
\end{Shaded}

But be aware that the statement-based binlog may cause inconsistency in replication, accroding to the warning in mysql manual:

\begin{quote}
When using statement-based logging for replication, it is possible for the data on the master and slave to become different if a statement is designed in such a way that the data modification is nondeterministic; that is, it is left to the will of the query optimizer. In general, this is not a good practice even outside of replication.
\end{quote}

So in general maybe it's better to use row-based logging, however that could result in larger binlog files.

\section{events in binlog}\label{events-in-binlog}

\subsection{View events in binlog}\label{view-events-in-binlog}

We can show the events contained in binlog use \texttt{show\ binlog\ events} command:

\begin{verbatim}
mysql> show binlog events in 'mysql-bin.000001';
+------------------+-----+----------------+-----------+-------------+---------------------------------------+
| Log_name         | Pos | Event_type     | Server_id | End_log_pos | Info                                  |
+------------------+-----+----------------+-----------+-------------+---------------------------------------+
| mysql-bin.000001 |   4 | Format_desc    |         1 |         123 | Server ver: 5.7.30-log, Binlog ver: 4 |
| mysql-bin.000001 | 123 | Previous_gtids |         1 |         154 |                                       |
+------------------+-----+----------------+-----------+-------------+---------------------------------------+
2 rows in set (0.00 sec)
\end{verbatim}

So let's try to find what happens when we do some operation in database. We'll create a new database and table, then insert a single record to the table by using the following commands:

\begin{Shaded}
\begin{Highlighting}[]
\KeywordTok{create} \KeywordTok{database}\NormalTok{ test;}
\KeywordTok{use}\NormalTok{ test;}
\KeywordTok{create} \KeywordTok{table}\NormalTok{ foo(}
    \KeywordTok{id} \DataTypeTok{int} \KeywordTok{not} \KeywordTok{null} \KeywordTok{primary} \KeywordTok{key}\NormalTok{,}
\NormalTok{    remark }\DataTypeTok{varchar}\NormalTok{(}\DecValTok{100}\NormalTok{)}
\NormalTok{);}
\KeywordTok{insert} \KeywordTok{into}\NormalTok{ foo(}\KeywordTok{id}\NormalTok{, remark) }\KeywordTok{values}\NormalTok{(}\DecValTok{1}\NormalTok{, }\StringTok{\textquotesingle{}hello world!\textquotesingle{}}\NormalTok{);}
\end{Highlighting}
\end{Shaded}

After this we'll find the following events has been appended to the binlog:

\begin{verbatim}
| mysql-bin.000001 | 219 | Query          |         1 |         313 | create database test                                                             |
| mysql-bin.000001 | 313 | Anonymous_Gtid |         1 |         378 | SET @@SESSION.GTID_NEXT= 'ANONYMOUS'                                             |
| mysql-bin.000001 | 378 | Query          |         1 |         520 | use `test`; create table foo(
id int not null primary key,
remark varchar(100)
) |
| mysql-bin.000001 | 520 | Anonymous_Gtid |         1 |         585 | SET @@SESSION.GTID_NEXT= 'ANONYMOUS'                                             |
| mysql-bin.000001 | 585 | Query          |         1 |         657 | BEGIN                                                                            |
| mysql-bin.000001 | 657 | Table_map      |         1 |         706 | table_id: 116 (test.foo)                                                         |
| mysql-bin.000001 | 706 | Write_rows     |         1 |         760 | table_id: 116 flags: STMT_END_F                                                  |
| mysql-bin.000001 | 760 | Xid            |         1 |         791 | COMMIT /* xid=26 */                                                              |
+------------------+-----+----------------+-----------+-------------+----------------------------------------------------------------------------------+
\end{verbatim}

\subsection{event types}\label{event-types}

As we already see, there are different kinds of events shown in binlog such as \texttt{Query} and \texttt{Table\_map}. Actually there are many more types than that, here are parts of them:

{\def\LTcaptype{none} % do not increment counter
\begin{longtable}[]{@{}ll@{}}
\toprule\noalign{}
Event type & Description \\
\midrule\noalign{}
\endhead
\bottomrule\noalign{}
\endlastfoot
UNKNOWN\_EVENT & never occurs \\
START\_EVENT\_V3 & a descriptor at each begging of binlog, and is replaced by FORMAT\_DESCRIPTION\_EVENT since MySQL 5.0 \\
QUERY\_EVENT & occurs when an updating statement is done \\
STOP\_EVENT & occurs when mysqld stops \\
INTVAR\_EVENT & occurs each time when a AUTO\_INCREMENT field is used \\
TABLE\_MAP\_EVENT & map the table defination to a number, happens before row operations \\
WRITE\_ROWS\_EVENT & insert records \\
UPDATE\_ROWS\_EVENT & update rows \\
DELETE\_ROWS\_EVENT & delete rows \\
\end{longtable}
}

For more detailed explaination about event types, please refer to \href{https://dev.mysql.com/doc/internals/en/event-meanings.html}{Event Meanings}.

References:

\begin{itemize}
\tightlist
\item
  \href{https://dev.mysql.com/doc/refman/5.7/en/replication-options.html}{MySQL 5.7 Reference Manual - Replication and Binary Logging Options and Variables}
\item
  \href{https://dev.mysql.com/doc/refman/5.7/en/binary-log.html}{MySQL 5.7 Reference Manual - The Binary Log}
\end{itemize}


\end{document}
